\section{Live-debug gotchas}

\begin{frame}{Five blockers we hit live (and the patches)}
  \footnotesize
  \begin{enumerate}
    \item \term{\cmd{DigitalDesign.testbench()} mandatory}
          (commit \cmd{e386cf2}). Pre-staged configs that predated
          this commit failed because \cmd{RtlSimRunner.run()}
          calls \cmd{design.testbench()}.
          \emph{Fix:} declare \cmd{EDA\_AGENTS\_TB\_DRIVER /
          TARGET / DIR / ENV} in \cmd{config.yaml}.
    \item \term{\cmd{CocotbDriver} requires \cmd{target=make sim}}.
          The driver dispatches on \cmd{target.startswith("make")};
          plain \cmd{target=sim} falls back to a \cmd{python3
          target} invocation that does not work for cocotb.
    \item \term{cocotb subprocess needs the librelane venv on
          \cmd{PATH}}. The Makefile embeds \cmd{cocotb-config
          --makefiles}; without
          \cmd{/path/to/librelane/.venv/bin} in \cmd{PATH} that
          shells out to nothing.
    \item \term{LibreLane v3 strict schema rejects
          \cmd{EDA\_AGENTS\_*} keys}.
          \emph{Fix today (uncommitted):}
          \cmd{LibreLaneRunner.\_write\_librelane\_config()} writes
          a sibling \cmd{config.librelane.yaml} with the namespace
          stripped; \cmd{dir::} paths still resolve.
    \item \term{GL sim glob mismatch}: LibreLane writes the
          synthesised netlist using raw \cmd{DESIGN\_NAME}
          (\cmd{demo\_goertzel\_fp32.nl.v}) but
          \cmd{GenericDesign.project\_name()} returns the
          kebab-cased \cmd{demo-goertzel-fp32}.
          \emph{Fix today (uncommitted):}
          \cmd{GlSimRunner.\_find\_post\_synth\_netlist} now tries
          both spellings plus a generic
          \cmd{*-yosys-synthesis/*.nl.v} fallback.
  \end{enumerate}
  \note{%
    Five real, small framework patches. Each diff is
    5-30 lines; cumulative shape: every one of these failures
    surfaced because the loop refused to silently skip a stage.
    That refusal is the design point -- we want gates to bark, not
    the loop to hide. Speakers should keep this slide short and
    skip ahead unless the audience asks for the details. The
    deeper takeaway: the abstractions held; nothing in the loop or
    the backend dispatch needed touching to fix any of these.}
\end{frame}
